% ---
% title: Two sales appointments
% tags: random-variables, independence, expectation
% ---

\begin{problem}
	A salesman has two appointments today. The first appointment leads to a
	sale with probability $0.3$. The second one, independently of the first,
	leads to a sale with probability $0.6$. Each sale is either of the deluxe
	model, which costs 1000 euros, or of the standard model, which costs
	500 euros. The two models are equally likely, independently of everything
	else. Let $X$ be the total value of today's sales (in euros).
	(a) Find the probability mass function of $X$.
	(b) Find $\E[X]$.
\end{problem}

\begin{solution}
	\idea First decide how many sales there are (0, 1 or 2) and at which
	appointments. Then, for each case, look at the models sold. Independence
	lets us multiply the probabilities along each case.
	\step{(a)} \emph{The possible values.} With no sale $X=0$. With one sale
	$X$ is $500$ or $1000$. With two sales the total is $500+500=1000$,
	$500+1000=1500$ or $1000+1000=2000$. So $X$ takes the values
	$0$, $500$, $1000$, $1500$, $2000$.

	\emph{The number of sales.} Let $S_1$ and $S_2$ be the events that the
	first and the second appointment lead to a sale. They are independent, and
	so are their complements. Hence
	\[ \begin{aligned}
		P(\text{no sale})            & =P(S_1^cS_2^c)=0.7\cdot0.4=0.28,  \\
		P(\text{only the 1st sells}) & =P(S_1S_2^c)=0.3\cdot0.4=0.12,    \\
		P(\text{only the 2nd sells}) & =P(S_1^cS_2)=0.7\cdot0.6=0.42,    \\
		P(\text{both sell})          & =P(S_1S_2)=0.3\cdot0.6=0.18.
	\end{aligned} \]
	These four cases are mutually exclusive and cover everything; indeed
	$0.28+0.12+0.42+0.18=1$.

	\emph{The models.} If there is exactly one sale, it is standard (500)
	or deluxe (1000) with probability $0.5$ each. If there are two sales, the
	two models are independent, so each of the four combinations
	(standard--standard, standard--deluxe, deluxe--standard, deluxe--deluxe)
	has probability $0.5\cdot0.5=0.25$.

	\emph{Putting it together.} For each value of $X$ we add the
	probabilities of the mutually exclusive ways it can happen.
	\begin{itemize}
		\item $X=0$: no sale, $p(0)=0.28$.
		\item $X=500$: exactly one sale, and it is standard:
		      \[ p(500)=(0.12+0.42)\cdot0.5=0.06+0.21=0.27. \]
		\item $X=1000$: either exactly one sale and it is deluxe, or two
		      sales and both are standard:
		      \[ p(1000)=(0.12+0.42)\cdot0.5+0.18\cdot0.25=0.27+0.045=0.315. \]
		\item $X=1500$: two sales, one standard and one deluxe (two orders):
		      \[ p(1500)=0.18\cdot(0.25+0.25)=0.18\cdot0.5=0.09. \]
		\item $X=2000$: two sales, both deluxe:
		      \[ p(2000)=0.18\cdot0.25=0.045. \]
	\end{itemize}
	In a table:
	\begin{center}
		\begin{tabular}{c|ccccc}
			$x$    & $0$    & $500$  & $1000$  & $1500$ & $2000$  \\ \hline
			$p(x)$ & $0.28$ & $0.27$ & $0.315$ & $0.09$ & $0.045$
		\end{tabular}
	\end{center}
	Check: $0.28+0.27+0.315+0.09+0.045=1$, as it must be for a probability
	mass function.
	\step{(b)} The expected value is the sum of the values, each weighted by
	its probability:
	\[
		\begin{aligned}
			\E[X] & =0\cdot0.28+500\cdot0.27+1000\cdot0.315+1500\cdot0.09+2000\cdot0.045 \\
			      & =0+135+315+135+90=675.
		\end{aligned}
	\]
	\begin{answer}
		(a) $p(0)=0.28$, $p(500)=0.27$, $p(1000)=0.315$, $p(1500)=0.09$,
		$p(2000)=0.045$ \quad (b) \mbox{$\E[X]=675$ euros}
	\end{answer}
	\begin{remark}
		\textbf{Common mistake:} forgetting that $X=1000$ can happen in two
		different ways (one deluxe sale, or two standard sales). Then the
		probabilities do not add up to $1$.

		\textbf{A shortcut for (b).} The value of one sale is $500$ or $1000$
		with probability $0.5$ each, so its expected value is
		$0.5\cdot500+0.5\cdot1000=750$ euros. The expected number of sales
		is $0.3+0.6=0.9$. Since the expected value of a sum is the sum of the
		expected values, $\E[X]=0.3\cdot750+0.6\cdot750=0.9\cdot750=675$
		euros, the same result.
	\end{remark}
\end{solution}
